% SOURCE_STATEMENT: "How can the interacting Green functions of many-body quantum systems be obtained from a functional differential equation rather than the usual Dyson integral equation?"
% Preparation and final self-audit: CONTRIBUTING.md prompts C and D; source/audit artifacts retained in the local physics_statements_2026_10_11 batch.
% !TeX program = lualatex
\documentclass[11pt]{article}
\usepackage[margin=25mm]{geometry}
\usepackage{fontspec}
\setmainfont{Latin Modern Roman}
\usepackage{amsmath,amssymb,mathtools}
\usepackage{unicode-math}
\setmathfont{Latin Modern Math}
\usepackage{xurl,hyperref}
\hypersetup{colorlinks=true,urlcolor=blue}
\setcounter{secnumdepth}{0}
\setlength{\parindent}{0pt}
\setlength{\parskip}{5pt}
\setlength{\emergencystretch}{3em}
\newcommand{\sref}[2]{\hyperlink{source:#1:#2}{[S#2]}}
\newcommand{\cataloguescope}[1]{\paragraph{Recorded target.}#1}
\newcommand{\R}{\mathbb R}
\newcommand{\C}{\mathbb C}
\newcommand{\N}{\mathbb N}
\newcommand{\Z}{\mathbb Z}
\newcommand{\Q}{\mathbb Q}
\newcommand{\F}{\mathbb F}
\newcommand{\eps}{\varepsilon}
\begin{document}
% UnsolvedMath ID 30003753; source catalogue code OWR-15998-003
\setcounter{section}{30003752}
\section{Functional Differential Equations for Many-Body Green Functions}\label{30003753}
{\small UnsolvedMath ID 30003753 · Mathematical Physics}
\subsection{Definitions and mathematical statement}

\paragraph{Shared notation.}$\mathbb N=\{1,2,\ldots\}$; the standard number fields have their usual meanings. All additional conventions are stated locally.


\paragraph{Definitions and conventions.}
An interacting one-body Green function is a time-ordered or otherwise conventionally ordered two-point correlation function of creation and annihilation fields. The Dyson equation writes $G=G_0+G_0\Sigma G$, with convolution in the space, spin and time variables and self-energy $\Sigma$. A functional differential equation treats $G$ as a functional of an added external potential and differentiates with respect to that potential, thereby encoding higher correlations without assuming a known self-energy.

\paragraph{Target retained from the canonical source catalogue.}
\paragraph{Catalogue wording retained for traceability.}
How can the interacting Green functions of many-body quantum systems be obtained from a functional differential equation rather than the usual Dyson integral equation?

\paragraph{Source scope and qualification.}
The source\textquotesingle s ordering, contour and external-field convention fix the functional equation. Producing formal perturbative derivatives alone does not construct the desired nonperturbative interacting Green function.
\subsection{Short English statement}
Obtain interacting many-body Green functions directly from their equations in functional derivatives.
\subsection{Sources}
\begin{itemize}
\item[\hypertarget{source:30003753:1}{S1}] ULAM.AI and the UnsolvedMath Contributors, UnsolvedMath, record 30003753 (OWR-15998-003). Catalogue curation under CC BY 4.0; definitions and English exposition added locally.
\url{https://huggingface.co/datasets/ulamai/UnsolvedMath}.
\item[\hypertarget{source:30003753:2}{S2}] Mathematical Methods in Quantum Chemistry (2018), Oberwolfach Reports, DOI 10.4171/owr/2018/13, pp. 631--721. Primary report accessed 2026-10-11; relevant PDF pages 56.
\url{https://doi.org/10.4171/owr/2018/13}.
\item[\hypertarget{source:30003753:3}{S3}] Bold Feynman diagrams (catalogue-listed context; not independently consulted in this preparation).
\url{https://arxiv.org/abs/1809.02901}.
\item[\hypertarget{source:30003753:4}{S4}] Coupled-cluster Green functions (catalogue-listed context; not independently consulted in this preparation).
\url{https://pmc.ncbi.nlm.nih.gov/articles/PMC8493088/}.
\item[\hypertarget{source:30003753:5}{S5}] Green-function implementation (catalogue-listed context; not independently consulted in this preparation).
\url{https://doi.org/10.1016/j.cpc.2024.109380}.
\end{itemize}
\label{end:30003753}
\end{document}
